\section{Предметная область междоменной маршрутизации и BGP}
\label{sec:bgp-domain}

\subsection{Автономные системы и междоменная маршрутизация}

Интернет представляет собой объединение большого числа независимых
сетей, принадлежащих различным организациям: операторам связи, облачным
платформам, корпорациям, университетам, государственным структурам и
другим участникам. На уровне административного управления эти сети
объединяются в автономные системы (Autonomous Systems, AS). Каждая
автономная система управляется единым субъектом и реализует собственную
политику маршрутизации.

Междоменная маршрутизация отличается от внутридоменной тем, что её цель
не ограничивается поиском кратчайшего пути. Здесь ключевую роль играют
политические и экономические ограничения: приоритет партнёрских связей,
экспортные и импортные политики, ограничения по префиксам и
соображения отказоустойчивости. Именно поэтому BGP, в отличие от
типичных IGP-протоколов, является policy-aware протоколом и работает не
как алгоритм минимизации расстояния, а как механизм обмена
достижимостью с учётом локальных политик
\cite{rfc4271,griffin2002stablepaths}.

\subsection{BGP как базовый протокол междоменной маршрутизации}

Согласно RFC~4271 \cite{rfc4271}, BGP-4 является междоменным протоколом
маршрутизации, основная задача которого состоит в обмене информацией о
достижимости сетевых префиксов между BGP-speaking системами.
Маршрутная информация в BGP включает не только факт достижимости
префикса, но и последовательность автономных систем (AS-path), через
которые эта достижимость проходит. Это позволяет одновременно выявлять
петли маршрутизации и принимать policy-based решения о выборе лучшего
маршрута.

Для целей настоящей работы важны несколько свойств BGP.

Во-первых, BGP оперирует префиксами, а не отдельными адресами.
Следовательно, задача сопоставления IP-адреса автономной системе в
конечном счёте сводится к longest prefix match по актуальному набору
объявленных префиксов.

Во-вторых, BGP изначально поддерживает classless routing и CIDR
\cite{rfc4632}: маршруты различной длины маски могут сосуществовать,
перекрываться и конкурировать за роль лучшего маршрута для конкретного
адреса.

В-третьих, BGP представляет состояние маршрутизации не как статическую
таблицу, а как непрерывно эволюционирующее множество объявлений
(announcements) и withdraw-событий. Это означает, что программная
система, желающая использовать BGP как источник истины, не может
ограничиться редким чтением снимка без риска работы с устаревшим
представлением.

\subsection{Multiprotocol extensions и поддержка нескольких семейств адресов}

RFC~4760 \cite{rfc4760} вводит расширения BGP-4 для переноса информации
о достижимости нескольких сетевых протоколов одновременно. Для этого
используются пары $\langle \mathit{AFI}, \mathit{SAFI} \rangle$
(Address Family Identifier и Subsequent Address Family Identifier) и
специальные атрибуты \texttt{MP\_REACH\_NLRI} и
\texttt{MP\_UNREACH\_NLRI}. Для рассматриваемой задачи это означает,
что прикладной сервис должен корректно работать как минимум с двумя
основными случаями:

\begin{itemize}
\item IPv4 unicast ($\mathit{AFI}=1$, $\mathit{SAFI}=1$);
\item IPv6 unicast ($\mathit{AFI}=2$, $\mathit{SAFI}=1$).
\end{itemize}

Различие между IPv4 и IPv6 важно не только на уровне протокола, но и на
уровне внутреннего поиска. Для IPv4 максимальная длина префикса
составляет 32 бита, для IPv6~-- 128~бит. Соответственно, реализация
longest prefix match, основанная на переборе возможных длин маски или
на байтовой навигации по структуре данных, имеет принципиально
различные характеристики на этих двух семействах адресов; разница
проявляется как в стоимости отдельного LPM-запроса, так и в общем числе
кандидатных длин маски, которые может потребоваться проверить.

\subsection{CIDR и longest prefix match как фундамент задачи}

RFC~4632 \cite{rfc4632} закрепляет практику Classless Inter-Domain
Routing и фиксирует переход от классовой адресации к префиксной как
ключевой механизм ограничения роста глобального маршрутизационного
состояния. В classless-модели значение имеет не «класс сети», а явная
пара \textit{префикс} и \textit{длина маски}.

Для прикладного сервиса это означает следующее. Если для одного и того
же адресного пространства одновременно объявлено несколько
перекрывающихся префиксов, то принадлежность конкретного IP-адреса
определяется по наиболее специфичному (longest) из них. Простого
словаря «адрес $\rightarrow$ значение» здесь принципиально
недостаточно. Внутреннее хранилище обязано либо нативно поддерживать
LPM, либо предоставлять структуру, поверх которой LPM может быть
эффективно реализован.

\subsection{Route Refresh и повторная синхронизация состояния}

RFC~2918 \cite{rfc2918} описывает механизм Route Refresh Capability,
позволяющий BGP-speaker запрашивать переобъявление соответствующего
\texttt{Adj-RIB-Out} у соседа без полного перезапуска BGP-сессии.
RFC~7313 \cite{rfc7313} расширяет этот механизм, вводя маркеры начала
и завершения refresh-последовательности (Beginning of RR, End of RR),
что позволяет более корректно и менее разрушительно проверять и
корректировать несогласованное состояние.

Для сервиса, работающего по схеме «снимок~+~поток обновлений», это
особенно важно. Потоковая модель почти неизбежно сталкивается с
вопросом ресинхронизации:

\begin{itemize}
\item что делать после разрыва соединения с источником данных;
\item как убедиться, что между двумя моментами наблюдения не были
потеряны withdraw-события;
\item как безопасно вернуть состояние к согласованному виду без полной
остановки и перезапуска всей системы.
\end{itemize}

Механизмы Route Refresh позволяют рассматривать такие сценарии не как
ad-hoc инженерный костыль, а как легитимную часть протокольной модели
BGP.

Дополнительный родственный сюжет -- BGP route flap damping
\cite{rfc2439}. Хотя данная работа не реализует damping явно, его
теоретический смысл, состоящий в подавлении нестабильности и
сглаживании церемоний объявления/отзыва, важен для понимания того,
почему обработка update-стрима не может быть наивной.

\subsection{BGP Monitoring Protocol и наблюдаемость состояния}

RFC~7854 \cite{rfc7854} определяет BGP Monitoring Protocol (BMP) как
специальный протокол для мониторинга BGP-сессий. В контексте настоящей
работы BMP интересен не как непосредственно используемый транспорт, а
как иллюстрация общей архитектурной идеи: маршрутизационное состояние
имеет смысл собирать, наблюдать и передавать как отдельный поток
данных, а не только как побочный эффект работы маршрутизатора.

Хотя в реализации, рассматриваемой в работе, интеграция строится не
через BMP, а через GoBGP gRPC API, концептуально оба подхода
принадлежат одному классу решений: отделение маршрутизационного plane
от сервисов, использующих его как источник данных.

%==============================================================
